\scanpage{12}

donc $\varphi\otimes\varphi'\in P(M\otimes M')$. Enfin, supposons $\varphi\in P(M)$, i.e. $\varphi\sim\varphi_{\underline m}|M$, alors $\varphi(i)\sim\varphi_{\underline m}(i)|M(i)\unsure{;|,}=\varphi_{(m_1,\ldots,m_r,m_{r+1}-i)}|M(i)$, donc $\varphi(i)\in P(M(i))$.

15) Soit $G$ groupe algébrique sur $\mathbb R$, \deleted{$C\in$} avec

\[
G_m\xrightarrow{i}G\xrightarrow{\varepsilon}G_m,\quad\unsure{y|\gamma|\eta}=i(-1),\quad\text{et soit }C\in G(\mathbb R)
\]

avec $C^2=y$. On prend $\mathcal M=\operatorname{Rep}(G)$

\deleted{1) Soit $M$} a) Si $M\in\operatorname{Ob}\mathcal M^i$ i.e. \inserted{muni d'une} repr de \deleted{degré} poids $i$, \deleted{\unsure{Obs.}} si $\varphi:M\otimes M\unsure{-|\to}\mathbb R(-i)$ est telle que $\varphi(x,Cy)$ est symétrique définie positive, alors $\varphi$ est une forme de Weil, et si on a de $=$ $\psi$ \deleted{\unsure{$\longmapsto$}} $:N\otimes N-\mathbb R(-i)$, $\psi(x,Cy)$ déf. positive, alors $\varphi$ et $\psi$ sont \uline{compatibles}.

\margin{$\varphi$ est dite une forme de $C$-polarisation \deleted{$\varphi$ est}}

% slanted
\margin{vérifions d'abord que $\varphi$ est \deleted{\unsure{sym}} $(-1)^i$-sym.}

En effet, soit $u:M\longrightarrow N$ compatible avec l'action de $G$ : Considérons $u^{*(\varphi,\psi)}$, et comparons avec $u^{*(\varphi^C,\psi^C)}$ ($\varphi^C(x,y)=\varphi(x,Cy)$). On aura

\[
\begin{aligned}
u^{*(\varphi^C,\psi^C)}&=C_M^{-1}u^{*(\varphi,\psi)}C_N\qquad\text{donc}\\
u^{*(\varphi,\psi)}&=C_M\underbrace{u^{*(\varphi^C,\psi^C)}}_{u'}C_N^{-1},\qquad\text{donc}\\
\operatorname{Tr}uu^{*(\varphi,\psi)}&=\operatorname{Tr}uC_Mu'C_N^{-1}
=\operatorname{Tr}\underbrace{C_N^{-1}u}_{uC_M^{-1}}C_Mu'=\operatorname{Tr}uu'\geq0
\end{aligned}
\]

% slanted
\margin{\unsure{dis q. tjrs la même|dis q. tjrs le même} (sans condition de positivité) $\operatorname{Tr}uu^{*(\varphi^{\unsure{*|C}},\psi)}\deleted{\unsure{0}}=\operatorname{Tr}uu^{*(\varphi^C,\psi^C)}$}

\deleted{Or $C_M$ et $C_N$ sont orthogonaux pour les formes déf positives $\varphi^C,\psi^C$ (en fait, les $g\in G(\mathbb R)$ qui dans $\operatorname{Centr}(C)$ opèrent de façon orthogonale), on a donc $C_N^{-1}=C$}

b) Si $\varphi,\psi$ sont des $C$-polarisations, $\varphi\otimes\psi$ aussi, si $\varphi$ une $C$-polarisation, alors $\varphi(i)$ aussi, \inserted{\uline{pourvu que} $\varepsilon(C)=1$ (et non $-1$ !)} la forme \unsure{standart|standard} $\mathbb R\otimes\mathbb R-\mathbb R$ est une $C$-polarisation.

% archive sheet number: 11
